





\usepackage{amsmath, amssymb, amsfonts}
\usepackage{geometry}
\usepackage{hyperref}
\usepackage{setspace}

% A4 sayfa kenar boşlukları
\geometry{
    a4paper,
    left=2.5cm,
    right=2.5cm,
    top=3cm,
    bottom=3cm
}

% Makale Başlığı ve Yazar Bilgileri
\title{\textbf{Bütünsel Kaynak Teorisi (UST) -- Versiyon 0}\\ 
\Large Master Axiom: Ontolojik Temel ve Omnium Sayı Doğrusu}
\author{
    \textbf{Niyazi ÖCAL} \\ 
    \textit{Bağımsız Araştırmacı / Bilgisayar Programcısı} \\
    \textit{Patent Başvuru No: 2026/003258} \\
    \href{mailto:niyazi.ocal@gmail.com}{niyazi.ocal@gmail.com}
}
\date{\today}

\begin{document}

\maketitle

\begin{abstract}
Bütünsel Kaynak Teorisi (UST) v0, Kuantum Mekaniği ve Genel Görelilik arasındaki uyuşmazlığı evrensel ve boyutsuz bir sabit olan $N_{s,q} \approx 0.63354460$ üzerinden çözen nihai ontolojik ve matematiksel çerçevedir. Bu çalışma, fiziksel evrenin temelini geleneksel parçacıklar yerine "Omnium" (Sıfırıncı Element) adı verilen pre-baryonik bir kök hücre olarak tanımlar. Geleneksel sayı doğrusu, merkezdeki Omnium çekirdeğinden dışarıya doğru yayılan çift yönlü bir "Omnium Sayı Doğrusu (ONL)" olarak yeniden yapılandırılarak; dondurulmuş kütleçekimsel arka plan (Kanal C) ile aktif kuantum sektörü (Kanal Q) arasında deterministik bir köprü kurulmuştur. Evrendeki bilgi bütünlüğü, donmuş potansiyel ile fiziksel tezahür arasında hiçbir bilgi kaybı yaşanmamasını sağlayan iki taraflı bir Limit Simetrisi ile güvence altına alınmıştır.
\end{abstract}

\vspace{1cm}

\section{Ontolojik Temel – Kök Hücre ve Omnium Sayı Doğrusu}

\subsection{Sıfırıncı Element: Omnium}
Bütünsel Kaynak Teorisi (UST), fiziksel evreni tanımlamaya geleneksel parçacık fiziği ile değil, temel bir "Kozmik Kök Hücre" (Source Cell) kavramı ile başlar. Bu kavram, periyodik tablonun ve fiziksel formların öncesinde yer alan, \textbf{Omnium (Sıfırıncı Element / Element Zero)}'dur. Klasik baryonik maddenin aksine Omnium; atom numarası, içsel yörüngesi, yükü veya spini olmayan, pre-baryonik (madde öncesi) saf bir potansiyel durumudur. Evrendeki tüm fiziksel özellikler, aslında bu saf potansiyelin $N_{s,q}$ operatörü tarafından işlenerek gözlemlenebilir boyutlara faz geçişi yaşamasıyla ortaya çıkan "ikincil" (emergent) karakteristiklerdir.

\subsection{Omnium Sayı Doğrusu (ONL) ve İki Kanallı Yapı}
Geleneksel matematik, sayı doğrusunu doğrusal bir ilerleyiş olarak kabul eder. UST ise sayısal sürekliliği (continuum), merkezdeki Omnium çekirdeğinden dışarıya doğru yayılan çift yönlü bir yapı olarak yeniden yapılandırır. Bu yapıya \textbf{Omnium Sayı Doğrusu (ONL)} denir:

\begin{equation}
\{ \dots, -n, \dots, -1, -0 \} \longleftarrow [\text{OMNİUM}] \longrightarrow \{ 0+, 1, \dots, n, \dots \}
\end{equation}

Bu yeni mimaride evren iki temel faza ayrılır:
\begin{itemize}
    \item \textbf{-0 (Giriş Fazı / Kanal C):} Bilginin statik bir iz (imprint) olarak depolandığı, zamanın akmadığı donmuş kütleçekimsel arka planı (Karanlık Madde) temsil eder.
    \item \textbf{0+ (Çıkış Fazı / Kanal Q):} Fiziksel tezahürün gerçekleştiği, deneyimlediğimiz aktif kuantum bilgi sektörünü temsil eder.
\end{itemize}
Bu iki faz (-0 ile 0+) arasındaki geçiş alanı ise Kuantum Tünellemesinin gerçekleştiği \textit{Omnium Köprüsü}'dür.

\subsection{Bilgi-Madde İkiliği: Limit Simetrisi}
Saf bilgi potansiyeli (-0) ile fiziksel tezahür (0+) arasındaki geçiş, rastgele değil, kusursuz bir simetri ile yönetilir. Bu geçiş, bilginin hiçbir şekilde kaybolmamasını sağlayan iki taraflı bir \textbf{Limit Simetrisi} ile formülize edilir:

\begin{equation}
\lim_{\epsilon \to 0^-} \Psi_{\text{Fiziksel}} = \Psi(-0) \longleftrightarrow \lim_{\epsilon \to 0^+} \Psi_{\text{Bilgi}} = \Psi(0+)
\end{equation}

Bu denklem, donmuş potansiyel durumu ile gözlemlenebilir durum arasındaki geçişte hiçbir bilgi kaybı yaşanmayacağını (üniter evrimin korunacağını) güvence altına alır. Madde ve bilgi, aslında bu simetrik sınırın iki farklı yüzünden ibarettir ve bu geçişi kontrol eden yegane şey evrensel algoritmanın çekirdeği olan $N_{s,q}$ sabitidir.

\end{document}